\section{Complete two-vertex dynamics}
\label{sec:direct-two-vertex}

Let $G$ be one edge with seed at vertex $0$. Both degrees equal one, so the
empirical score is simply $|r_u|/2$ and every batch contains one coordinate.
It is convenient to use the equivalent degree-scaled residual
\begin{equation}
 \widehat r:=\frac{2}{1+\alpha}D^{1/2}r.
 \label{eq:direct-scaled-residual}
\end{equation}
Set
\begin{equation}
 \gamma_\alpha:=\frac{2\alpha}{1+\alpha},
 \qquad
 c_\alpha:=\frac{1-\alpha}{1+\alpha}
 =\frac{2\lambda_\alpha}{1+\lambda_\alpha^2}.
 \label{eq:direct-scaled-parameters}
\end{equation}
The scaled system starts from
$\widehat r^{(0)}=(\gamma_\alpha,0)$; the terminal gate is
$\|\widehat r\|_\infty\leq\gamma_\alpha\epsppr$. An optimal-SOR push leaves
$-\lambda_\alpha^2\sigma$ at the pushed coordinate and sends
$2\lambda_\alpha\sigma$ to its neighbor.

\begin{lemma}[Closed-form ranked trajectory on one edge]
\label{lem:direct-edge-trajectory}
Write $\lambda:=\lambda_\alpha$. The score-maximizing spreading phase
alternates coordinates. After $2k$ and $2k+1$ pushes, respectively,
\begin{align}
 \widehat r^{(2k)}
 &=\gamma_\alpha\lambda^{2k}
   \bigl(2k+1,-2k\lambda\bigr),
 \label{eq:direct-even-trajectory}\\
 \widehat r^{(2k+1)}
 &=\gamma_\alpha\lambda^{2k}
   \bigl(-\lambda^2(2k+1),2\lambda(k+1)\bigr).
 \label{eq:direct-odd-trajectory}
\end{align}
The first coordinate has larger magnitude in every even state, and the
second has larger magnitude in every odd state.
\end{lemma}

\begin{proof}
The claim holds at $k=0$. In an even state,
$2k+1>2k\lambda$, so coordinate $0$ is selected. Applying the optimal-SOR
update gives \cref{eq:direct-odd-trajectory}. There,
\begin{equation*}
 2\lambda(k+1)>\lambda^2(2k+1),
\end{equation*}
so coordinate $1$ is selected. One more update gives
\cref{eq:direct-even-trajectory} with $k$ replaced by $k+1$.
\end{proof}

The formulas exhibit critical damping: the geometric factor
$\lambda^{2k}$ is multiplied by the linear transient $k$. More importantly,
the two residual entries have opposite signs. The exact rung can exploit
that correlation.

\begin{lemma}[One-push settlement ratio]
\label{lem:direct-edge-settlement-ratio}
Suppose the spreading phase stops in state $2k$ or $2k+1$ and the exact phase
pushes the larger-magnitude coordinate. The remaining residual magnitude,
divided by the pre-push maximum magnitude, equals
\begin{align}
 R_{2k}(\lambda)
 &=\frac{2\lambda\{k(1-\lambda^2)+1\}}
         {(1+\lambda^2)(2k+1)},
 \label{eq:direct-even-settlement}\\
 R_{2k+1}(\lambda)
 &=\frac{\lambda\{2k(1-\lambda^2)+3-\lambda^2\}}
         {2(k+1)(1+\lambda^2)}.
 \label{eq:direct-odd-settlement}
\end{align}
Both ratios satisfy
\begin{equation}
 R_j(\lambda)
 \leq
 \frac{\lambda(1-\lambda^2)}{1+\lambda^2}
 +\frac1{k+1}.
 \label{eq:direct-settlement-upper}
\end{equation}
Furthermore,
\begin{equation}
 R_\star
 :=\max_{0<\lambda<1}
   \frac{\lambda(1-\lambda^2)}{1+\lambda^2}
 =0.3002831060\ldots,
 \label{eq:direct-universal-settlement}
\end{equation}
where the maximum occurs at
$\lambda^2=\sqrt5-2$.
\end{lemma}

\begin{proof}
An exact push sends $c_\alpha$ times its residual to the neighbor. Apply this
to \cref{eq:direct-even-trajectory,eq:direct-odd-trajectory}, use
$c_\alpha=2\lambda/(1+\lambda^2)$, and simplify to obtain
\cref{eq:direct-even-settlement,eq:direct-odd-settlement}.

For the even ratio, separate the term with $k(1-\lambda^2)$; the remaining
coefficient is at most $1/(2k+1)$. For the odd ratio, the remaining
coefficient is at most $1/(k+1)$ because
\begin{equation*}
 \frac{\lambda(3-\lambda^2)}{2(1+\lambda^2)}\leq1.
\end{equation*}
This proves \cref{eq:direct-settlement-upper}. Differentiating the limiting
function gives numerator
$1-4\lambda^2-\lambda^4$; its unique root in $(0,1)$ satisfies
$\lambda^2=\sqrt5-2$, proving \cref{eq:direct-universal-settlement}.
\end{proof}

\begin{theorem}[Eventual one-push terminal rung on one edge]
\label{thm:direct-edge-one-push}
Let
\begin{equation}
 B_{\rm edge}:=R_\star^{-1}=3.3301906768\ldots.
 \label{eq:direct-edge-band-ceiling}
\end{equation}
For every fixed $1<B<B_{\rm edge}$, there is an integer $K(B)$ such that if
the spreading phase executes at least $2K(B)$ pushes, then its terminal exact
phase performs at most one push. For $B=2.5$, one may take $K(B)=10$.
\end{theorem}

\begin{proof}
At the switch, the maximum scaled residual is at most
$B\gamma_\alpha\epsppr$. Choose $K(B)$ so that
\begin{equation*}
 B\left(R_\star+\frac1{K(B)+1}\right)\leq1.
\end{equation*}
After one exact push, \cref{lem:direct-edge-settlement-ratio} bounds the only
remaining residual by $\gamma_\alpha\epsppr$, while the pushed coordinate is
zero. Thus the terminal gate holds. For $B=2.5$,
\begin{equation*}
 2.5\left(0.3002831061+\frac1{11}\right)<1,
\end{equation*}
so $K=10$ suffices.
\end{proof}

\begin{corollary}[Accelerated one-edge work]
\label{cor:direct-edge-work}
For $0<\alpha\leq1/9$, the number $N$ of spreading pushes is characterized,
up to absolute constant factors, by
\begin{equation}
 (k+1)\lambda_\alpha^{2k}\asymp B\epsppr,
 \qquad k=\left\lfloor\frac N2\right\rfloor.
 \label{eq:direct-edge-count-characterization}
\end{equation}
Consequently,
\begin{equation}
 N
 =\Theta\!\left(
   \frac{1+\log(1/(\epsppr\sqrt\alpha))}{\sqrt\alpha}
  \right)
 \label{eq:direct-edge-accelerated-count}
\end{equation}
in the joint small-$\alpha$, small-gate regime, up to the harmless iterated
logarithm from inverting $k\lambda_\alpha^{2k}$. The exact phase adds at most
one push once \cref{thm:direct-edge-one-push} applies.
\end{corollary}

\begin{proof}
For $\alpha\leq1/9$, one has $\lambda_\alpha\geq1/2$. The larger magnitude
in either \cref{eq:direct-even-trajectory,eq:direct-odd-trajectory}, divided
by $\gamma_\alpha$, lies between
$(k+1)\lambda_\alpha^{2k}$ and
$2(k+1)\lambda_\alpha^{2k}$. This proves
\cref{eq:direct-edge-count-characterization}. Finally,
$-\log\lambda_\alpha=2\operatorname{arctanh}(\sqrt\alpha)
=\Theta(\sqrt\alpha)$; invert the displayed geometric-polynomial relation.
\end{proof}

\begin{remark}[Why the self-reflection predictor is conservative]
The self-reflection-only condition would require
$B\lambda_\alpha^2\leq1$, or $B\leq\lambda_\alpha^{-2}$, which tends to one
as $\alpha\downarrow0$. The two-vertex theorem instead permits every fixed
$B<B_{\rm edge}$. The difference is exactly the opposite-sign neighbor
residual already present at the switch: the exact push combines with it and
cancels most of both packets. This is a rigorous mechanism by which the
measured $B=2.5$ can outperform a band rule derived from self-reflection
alone.
\end{remark}
